\documentclass[11pt]{article}

\usepackage[left=0.8in,right=0.8in,top=0.75in,bottom=0.75in]{geometry}
\usepackage[T1]{fontenc}
\usepackage[utf8]{inputenc}
\usepackage[none]{hyphenat}
\usepackage{setspace}

\pagestyle{empty}
\setlength{\parskip}{0.7em}
\setlength{\parindent}{0pt}

\begin{document}

\begin{flushleft}
Abin Daniel Zorto\\
Intelligent Technologies Research Group\\
University of East London\\
University Way, London E16 2RD, UK\\
\texttt{u2091940@uel.ac.uk}
\end{flushleft}

\vspace{0.8em}

\begin{flushleft}
Professor Panicos A. Kyriacou\\
Editor-in-Chief\\
\emph{Biomedical Signal Processing and Control}\\
Elsevier
\end{flushleft}

\vspace{0.8em}

\begin{flushleft}
\today
\end{flushleft}

\vspace{0.8em}

Dear Professor Kyriacou,

We wish to submit our original research article, \emph{Sparse Biomarker Discovery from Low-Montage Resting-State EEG for tDCS Treatment Response in Adults with Treatment-Resistant Depression}, for consideration for publication in \emph{Biomedical Signal Processing and Control}.

This study addresses patient-level prediction of treatment response from low-montage resting-state EEG in adults with treatment-resistant depression. Using pretreatment EEG from a fully remote sham-controlled tDCS trial, we evaluated a controlled sweep of window sizes and sparse feature-selection budgets across a hybrid deep model and a linear baseline under leave-one-participant-out validation, with class rebalancing and feature selection restricted to the training participants within each fold.

The main finding is that the strongest discriminatory performance emerged under an extremely sparse feature budget. The best hybrid configuration reached a patient-level ROC-AUC of 0.816, balanced accuracy of 0.786, and MCC of 0.571 across 21 held-out participants while selecting only one feature per fold and seven recurrent features overall. The most recurrent biomarkers centered on temporal spectral entropy and frontal-temporal difference measures, which preserved an interpretable signal profile alongside useful predictive performance.

We believe the article is well suited to \emph{Biomedical Signal Processing and Control} because it sits directly at the interface of biomedical signal analysis and clinically motivated decision support. The study develops and evaluates a practical EEG signal-processing and modeling framework for treatment-response stratification, emphasizes leakage-safe participant-level validation, and focuses on a low-montage setting with clear relevance for clinically deployable and remotely collected neurophysiological data. This fits the journal's emphasis on applications-led work in the measurement and analysis of biomedical signals for diagnosis, monitoring, and management.

This submission is original, has not been published previously, is not under consideration elsewhere, and has been approved by all authors.

Thank you for your time and consideration.

Sincerely,

\vspace{1.2em}

Abin Daniel Zorto\\
Corresponding author\\
\texttt{u2091940@uel.ac.uk}

\end{document}
